% ---------------------------------------------------------------------------------------
% obtune -- "Does fine-tuning teach semantic invariance?" -- FSE research-track draft v0.1
%
% 2026-09-12. STATUS: SKELETON. Argument is committed; several claims are still \slot.
%
% HOUSE RULE, inherited from paper_modularity/main.tex: NO NUMBER IS TYPED INTO THIS FILE.
% Every table under tables/ is emitted by scripts/paper/gen_tables.py from a JSON under
% results/analysis/, which is itself computed from the per-cell parquets in results/cells/.
% MASTER_REPORT.md and docs/ are an index of where to look, NOT a source; where a report and
% a result file disagree, the file wins. Regenerate with:
%     python scripts/paper/gen_tables.py
%
% TWO MACROS CARRY THE DRAFT STATE:
%   \pending -- a result that does not exist yet. NEVER replace with a plausible value.
%   \slot    -- a place where the argument needs a claim we have not earned yet.
% `grep -c 'pending\|slot' main.tex tables/*.tex` is the submission checklist; both reach zero.
%
% THREE THINGS THIS PAPER MUST NOT CLAIM, enforced by review against docs/PAPER_DRAFT.md 12:
%   1. that merging fails through task-vector interference  -- refuted, Table \ref{tab:geom}
%   2. that anchoring pays no clean-code cost               -- 5 of 8 models, Table \ref{tab:panel}
%   3. anything causal about model lineage                  -- n=8, p=0.07 at best, confounded
% ---------------------------------------------------------------------------------------

\documentclass[sigconf,review,anonymous]{acmart}

\settopmatter{printacmref=false}
\renewcommand\footnotetextcopyrightpermission[1]{}
\pagestyle{plain}
\acmConference[FSE '27]{Under review}{}{}
\setcopyright{none}

% acmart already loads booktabs, amsmath and newtxmath; adding amssymb here collides with
% newtxmath on \Bbbk and halts the engine. Do not re-add it.
\usepackage{xcolor}

% --- draft-state macros ---------------------------------------------------------------
\newcommand{\pending}{\textcolor{gray}{\textit{--}}}
\newcommand{\slot}[1]{\textbf{[CLAIM SLOT: #1]}}

% --- verdict marks used inside the generated tables ------------------------------------
% NB: these are \vok/\vrefuted/\vinc, NOT \ok/\ref/\inc. `\ref` is LaTeX's cross-reference
% command; redefining it makes every \ref{tab:...} in the prose render as a verdict mark.
\newcommand{\vok}{\textsuperscript{\checkmark}}
\newcommand{\vrefuted}{\textsuperscript{\texttimes}}
\newcommand{\vinc}{\textsuperscript{$\circ$}}

\begin{document}

\title{Does fine-tuning teach semantic invariance?\\
       Composition, anchoring, and what adaptation to obfuscated code actually learns}

\author{Anonymous Author(s)}
\affiliation{\institution{Under review}\country{}}

\begin{abstract}
Fine-tuning a code model on obfuscated programs raises its accuracy on those programs. We ask
whether that buys \emph{invariance} to meaning-preserving transformation or merely
\emph{memorisation} of the transforms seen in training, and we separate the two by never training
or evaluating on recovery: every evaluation is on code that stays obfuscated.

Across \textbf{eight models spanning four lineages} we find a dissociation that single-lineage
evaluation hides. Training on a \emph{breadth} of transforms composes on stacks built from those
same transforms, and pays a measurable tax on a held-out transform family --- a tax whose sign is
the same on every model we tested, with no contradiction. A paired-consistency objective that
anchors the model's answer distribution to its own behaviour on the unobfuscated parent removes
that tax on six of eight models, and \emph{reverses} on one. Its clean-code cost is likewise
model-dependent, holding on five of eight.

We then build the stimulus the question needs and that the literature lacks: stacks that contain a
transform family the model has never seen. Breadth's advantage flips sign the moment such a
component enters, while anchoring rises above the clean-code control. That flip is not uniform:
it is concentrated in stacks that leave the model no identifier surface to key on, which we
establish with a rule fixed before the stimulus existed. \slot{one-sentence takeaway once the
out-of-sample divergence reads complete}
\end{abstract}

\maketitle

% =======================================================================================
\section{Introduction}
\slot{the five moves, from docs/PAPER_DRAFT.md 3}

The paper's separation from the DOBF lineage is a design decision rather than a framing: we never
train or evaluate on deobfuscation. Evaluation is always output prediction over code that remains
obfuscated, so what we measure is semantic competence \emph{under} transformation and not the
ability to undo it.

% =======================================================================================
\section{Setup}
\label{sec:setup}

\paragraph{Condition ladder.} Conditions are \textbf{single-transform from the clean parent, never
stacked}, with identical semantics in Python and JavaScript: identifier-level (adversarial
renaming, random renaming, sequential minification) and structural (control-flow flattening, opaque
predicates with dead code). A held-out family --- string encoding plus guarded mixed
boolean-arithmetic rewriting --- is never trained on. Because its evaluation budget is spent, every
held-out number in this paper is measured on \textbf{X1}, a trainable sibling built from a
different surface scheme that reproduces the held-out family's difficulty level
($r=0.9992$).\footnote{The sibling's calibration and the quarantine discipline behind it are in the
supplement; no number here reads the quarantined family.}

\paragraph{Systems.} \textsc{clean} tunes on the unobfuscated parent only. \textsc{mono} tunes on
the breadth of all six ladder conditions. \textsc{cons} adds a paired-consistency term,
\[
\mathcal{L} = \mathrm{CE}(y \mid x_{\text{obf}})
  + \lambda\,\mathrm{KL}\!\left(p_{T}(\cdot \mid x_{\text{parent}}) \,\|\, p_{S}(\cdot \mid x_{\text{obf}})\right)
\]
at answer tokens, where $T$ is a \emph{frozen clean-code-tuned} adapter reading the unobfuscated
parent of the same program and case. \textsc{fam} tunes on the held-out family's sibling.
\textbf{\textsc{mono} is the exact data-matched control for \textsc{cons}}: both train on the same
26{,}841 rows for the same 1{,}257 steps, and with $\lambda=0$ the objective reduces to
\textsc{mono} identically --- the student never sees the parent view, only the frozen teacher does.

\paragraph{Statistics.} Every contrast is a program-clustered bootstrap (2{,}000 resamples, seed
17); programs rather than items are the resampling unit because one program contributes several
correlated input cases. Intervals are 95\%. Coverage differs by condition, so every contrast runs
on a common subset of programs \textbf{recorded before any system is evaluated}.

% =======================================================================================
\section{The measurement, before the results}
\label{sec:grid}

All eight models are read from a single evaluation phase with one configuration. That configuration
is newer than the four models published earlier, so before reading it we recompute each of those
models' previously published contrasts through it and ask whether the published point falls
\emph{inside} the new interval --- whether the two are the same measurement, not whether they round
alike.

\begin{table}[t]\centering\small
\caption{The uniform grid reproduces every previously published value it was checked against, so a
departure in Table~\ref{tab:panel} is a property of the model and not of the configuration.}
\label{tab:grid}
% GENERATED by scripts/paper/gen_tables.py on 2026-09-12 17:55 UTC
% SOURCE: results/analysis/pipeline/grid_agreement_panel_core.json
% Do not edit by hand -- regenerate. Numbers come from result files, never from a report.
\begin{tabular}{lcccc}
\toprule
model & R1 pub.\ $\rightarrow$ grid & R2 pub.\ $\rightarrow$ grid & R3 pub.\ $\rightarrow$ grid & max $|\Delta|$ \\
\midrule
CodeLlama-7B & $-3.79 \rightarrow -3.71$ & $+4.59 \rightarrow +5.44$ & $-0.30 \rightarrow -0.54$ & 0.85 \\
CodeLlama-13B & $-4.28 \rightarrow -4.04$ & $+3.95 \rightarrow +3.87$ & $-0.60 \rightarrow -0.72$ & 0.24 \\
CodeLlama-34B & $-2.64 \rightarrow -2.72$ & $+3.38 \rightarrow +3.62$ & $+0.42 \rightarrow +0.48$ & 0.24 \\
Llama-3.1-8B & $-2.88 \rightarrow -2.72$ & $+3.46 \rightarrow +3.38$ & $+0.00 \rightarrow -0.30$ & 0.30 \\
\bottomrule
\end{tabular}
\end{table}

% =======================================================================================
\section{RQ1 --- Does single-transform adaptation compose?}
\label{sec:rq1}

\paragraph{Composition in the data works, and only within what the data covered.} Breadth training
gains on stacks built from transforms it saw and does not on stacks containing a family it did not.
We state that as a \emph{difference} rather than as two levels, because one interval being positive
while another is negative is not a test of the difference between them --- and on most models here
the unseen-side interval straddles zero on its own. The two halves are paired on the programs common
to every cell of both groups, so the difference carries an interval.

\begin{table}[t]\centering\small
\caption{RQ1's dissociation, within model, on one program set. The \emph{difference} is the
measured quantity; the levels are context. StarCoder2 is the clearest case: both of its levels
straddle zero and its difference does not.}
\label{tab:dissoc}
% GENERATED by scripts/paper/gen_tables.py on 2026-09-12 17:55 UTC
% SOURCE: results/analysis/pipeline/stack_dissociation.json
% Do not edit by hand -- regenerate. Numbers come from result files, never from a report.
\begin{tabular}{llccc}
\toprule
model & lineage & seen stacks & unseen inside & \textbf{difference} \\
\midrule
CodeLlama-7B & Meta & \textbf{$+3.36$ \tiny[+1.57, +5.27]} & \textbf{$-1.92$ \tiny[-3.70, -0.14]} & \textbf{$+5.28$ \tiny[+3.20, +7.54]} \\
CodeLlama-13B & Meta & \pending & \pending & \pending \\
CodeLlama-34B & Meta & \textbf{$+2.27$ \tiny[+0.44, +4.09]} & $-0.77$ \tiny[-2.62, +1.12] & \textbf{$+3.03$ \tiny[+0.78, +5.28]} \\
Llama-3.1-8B & Meta & \textbf{$+2.30$ \tiny[+0.47, +4.15]} & $-1.05$ \tiny[-2.72, +0.63] & \textbf{$+3.35$ \tiny[+1.50, +5.30]} \\
StarCoder2-15B & BigCode & $+1.50$ \tiny[-0.02, +3.00] & $-1.19$ \tiny[-2.93, +0.49] & \textbf{$+2.68$ \tiny[+0.77, +4.65]} \\
Gemma-3-12B & Google & \pending & \pending & \pending \\
CodeGemma-7B & Google & \pending & \pending & \pending \\
Granite-3.1-8B & IBM & \pending & \pending & \pending \\
\midrule
\multicolumn{4}{l}{\emph{difference significant on}} & 4 of 4 read \\
\bottomrule
\end{tabular}
\end{table}

\slot{routing and merging results on stacks, from the modularity thread}

Weight-space merging fails, and we report that failure \textbf{without} the mechanism usually
offered for it. The adapters a merge is built from are the \emph{most} aligned and \emph{least}
conflicting bank we measured, and merging them still loses to them --- a mechanism predicting
cancellation cannot explain a failure that occurs where cancellation is lowest. What dominates the
geometry is not the task but the seed.

\begin{table}[t]\centering\small
\caption{Task-vector geometry. Varying the \emph{seed} alone moves the update direction far more
than varying the \emph{condition} does.}
\label{tab:geom}
% GENERATED by scripts/paper/gen_tables.py on 2026-09-12 18:11 UTC
% SOURCE: results/analysis/f1b/geometry_codellama7b_*.json
% Do not edit by hand -- regenerate. Numbers come from result files, never from a report.
\begin{tabular}{llccc}
\toprule
adapter bank & what varies & mean cosine & sign conflict & TIES kept \\
\midrule
5 specialists, one seed & condition only & 0.653 & 0.323 & 0.916 \\
clean-code arm, 3 seeds & seed only & 0.119 & 0.470 & 0.785 \\
5 specialists $\times$ 2 seeds & both & 0.344 & 0.554 & 0.699 \\
\bottomrule
\end{tabular}
\end{table}

% =======================================================================================
\section{RQ2 --- What is learned, and why it does not compose}
\label{sec:rq2}

\slot{the five instruments, from docs/PAPER_DRAFT.md 6}

Breadth's advantage on a stack depends on the stack containing an \textbf{identifier} transform.
We establish this as a difference rather than as an ordering, because one interval excluding zero
while another does not is not a test of the difference between them.

\begin{table}[t]\centering\small
\caption{Breadth's gain over the clean-code control, per composite, on stacks containing the
held-out family; and the identifier-versus-structural difference on two independent stimulus sets.
The rule defining the split was fixed \textbf{before} the unseen-containing stimulus existed.}
\label{tab:split}
% GENERATED by scripts/paper/gen_tables.py on 2026-09-12 14:11 UTC
% SOURCE: results/analysis/pipeline/{f2_divergence_codellama7b,f3a_stack_identifier_*}.json
% Do not edit by hand -- regenerate. Numbers come from result files, never from a report.
\begin{tabular}{llc}
\toprule
composite & composition & \textsc{mono}$-$\textsc{clean} \\
\midrule
\texttt{C\_L1r\_X1} & ident.\ $\to$ unseen & \textbf{$+2.67$ \tiny[+0.35, +5.11]} \\
\texttt{C\_L1r\_X1m} & ident.\ $\to$ unseen (MBA half) & $+1.28$ \tiny[-1.16, +3.72] \\
\texttt{C3\_L1r\_S1\_X1} & ident.\ $\to$ struct.\ $\to$ unseen & $-0.35$ \tiny[-2.67, +1.86] \\
\texttt{C\_S2\_X1} & struct.\ $\to$ unseen & \textbf{$-3.60$ \tiny[-6.40, -0.58]} \\
\texttt{C\_S1\_X1s} & struct.\ $\to$ unseen (str.\ half) & \textbf{$-3.84$ \tiny[-6.75, -0.93]} \\
\texttt{C\_X1\_S1} & unseen $\to$ struct. & \textbf{$-4.19$ \tiny[-7.09, -1.28]} \\
\midrule
\multicolumn{2}{l}{\emph{stacks containing the unseen family}: identifier $-$ structural} & \textbf{$+5.85$ \tiny[+3.10, +8.72]} \\
\multicolumn{2}{l}{\emph{all-seen depth-3/4 stacks}: identifier $-$ structural} & \textbf{$+3.39$ \tiny[+0.82, +6.16]} \\
\bottomrule
\end{tabular}
\end{table}

\paragraph{What we do not claim.} That structural transforms contribute nothing: the established
quantity is the \emph{difference} between the groups. And that what transfers within a family is
its surface scaffolding --- the test that discriminates the surface reading from a weakest-link
reading returned \textsc{undecided} on every one of its three pre-registered rules.

% =======================================================================================
\section{RQ3 --- Anchoring to clean-code behaviour}
\label{sec:rq3}

\begin{table*}[t]\centering\small
\caption{Eight models, four lineages, one evaluation phase. \vok\ replicated, \vinc\ inconclusive,
\vrefuted\ refuted. Bold intervals exclude zero.}
\label{tab:panel}
% GENERATED by scripts/paper/gen_tables.py on 2026-09-12 18:11 UTC
% SOURCE: results/analysis/panel_replication_final.json
% Do not edit by hand -- regenerate. Numbers come from result files, never from a report.
\begin{tabular}{llccc}
\toprule
model & lineage & \textbf{R1} $\Delta_{\mathrm{breadth}}$ & \textbf{R2} $\Delta_{\mathrm{anchor}}$ & \textbf{R3} $\Delta_{\mathrm{clean}}$ \\
 & & \textsc{mono}$-$\textsc{clean} @X1 & \textsc{cons}$-$\textsc{mono} @X1 & \textsc{cons}$-$\textsc{clean} @L0 \\
\midrule
CodeLlama-7B & Meta & \textbf{$-3.71$ \tiny[-5.93, -1.56]} \vok & \textbf{$+5.44$ \tiny[+3.37, +7.41]} \vok & $-0.54$ \tiny[-1.98, +1.08] \vok \\
CodeLlama-13B & Meta & \textbf{$-4.04$ \tiny[-6.26, -1.89]} \vok & \textbf{$+3.87$ \tiny[+2.06, +5.77]} \vok & $-0.72$ \tiny[-2.28, +0.90] \vok \\
CodeLlama-34B & Meta & \textbf{$-2.72$ \tiny[-5.02, -0.49]} \vok & \textbf{$+3.62$ \tiny[+1.40, +5.76]} \vok & $+0.48$ \tiny[-0.84, +1.80] \vok \\
Llama-3.1-8B & Meta & \textbf{$-2.72$ \tiny[-4.70, -0.91]} \vok & \textbf{$+3.38$ \tiny[+1.57, +5.36]} \vok & $-0.30$ \tiny[-1.98, +1.26] \vok \\
StarCoder2-15B & BigCode & \textbf{$-2.72$ \tiny[-4.94, -0.58]} \vok & \textbf{$+2.31$ \tiny[+0.58, +4.03]} \vok & $-0.96$ \tiny[-2.51, +0.60] \vok \\
Gemma-3-12B & Google & \textbf{$-4.70$ \tiny[-6.93, -2.55]} \vok & \textbf{$+2.64$ \tiny[+0.91, +4.53]} \vok & \textbf{$-1.80$ \tiny[-3.53, -0.06]} \vrefuted \\
CodeGemma-7B & Google & $-1.98$ \tiny[-3.95, +0.16] \vinc & $-0.41$ \tiny[-2.39, +1.49] \vinc & \textbf{$-2.04$ \tiny[-3.77, -0.18]} \vrefuted \\
Granite-3.1-8B & IBM & $-0.91$ \tiny[-2.97, +1.15] \vinc & \textbf{$-1.89$ \tiny[-3.54, -0.25]} \vrefuted & \textbf{$-3.41$ \tiny[-5.39, -1.26]} \vrefuted \\
\midrule
\multicolumn{2}{l}{\emph{replicated / inconclusive / refuted}} & 6/2/0 & 6/1/1 & 5/0/3 \\
\bottomrule
\end{tabular}
\end{table*}

Three readings, at the strength the panel supports.

\paragraph{The breadth tax is the robust one.} Its sign is the same on \emph{every} model, six
significantly and none contradicting, across four lineages and 7B--34B.

\paragraph{Anchoring's repair is conditional.} It holds on six models and \textbf{reverses} on one,
where the anchored arm is worse than breadth on the held-out family. We name that model rather than
relegating it.

\paragraph{The clean-code claim is conditional too.} ``Anchoring pays no clean-code cost'' holds on
five of eight. It was discovered on four models of one lineage and holds on all four of them.

\paragraph{A warning, not a result.} Every failure above falls outside the lineage the findings were
discovered on, and all three are four-of-four inside it. We report this as a \emph{methodological}
warning and not as a claim about lineages: Fisher's exact on the sharpest split gives $p=0.0714$
one-sided with $n=8$ models, and the two models of one vendor both account for failures of the same
finding, so vendor and lineage are perfectly confounded. The strongest non-lineage model in the
panel is also the one an untuned capability screen scored at zero and would have discarded ---
removing it would have made this pattern look \emph{cleaner} and be \emph{more wrong}.

% =======================================================================================
\section{RQ4 --- Robustness, or a better surface heuristic?}
\label{sec:rq4}

Every composite stimulus in the prior literature, and in our own earlier work, is built from
transforms the model has seen. The question RQ4 asks cannot be answered with such a stimulus. We
build six composites that each contain a component from the held-out family, fix the program subset
and five decision rules in advance, and commit \textbf{both} readings of the outcome before running.

\begin{table}[t]\centering\small
\caption{The divergence ladder. Breadth's gain flips sign once the stack contains a family the
model has never seen; anchoring rises above the clean-code control.}
\label{tab:ladder}
% GENERATED by scripts/paper/gen_tables.py on 2026-09-21 02:05 UTC
% SOURCE: results/analysis/pipeline/{f2_divergence_codellama7b,composite_depth_codellama7b}.json
% Do not edit by hand -- regenerate. Numbers come from result files, never from a report.
\begin{tabular}{llccc}
\toprule
level & stimulus & \textsc{mono}$-$\textsc{clean} & \textsc{cons}$-$\textsc{clean} & \textsc{cons}$-$\textsc{mono} \\
\midrule
d0 & depth 2, all seen & \textbf{$+3.49$ \tiny[+2.04, +5.01]} & \textbf{$+4.76$ \tiny[+3.53, +6.05]} & $+1.27$ \tiny[-0.02, +2.59] \\
d1 & depth 3--4, all seen & \textbf{$+4.33$ \tiny[+2.63, +6.03]} & \textbf{$+5.28$ \tiny[+3.79, +6.84]} & $+0.95$ \tiny[-0.37, +2.31] \\
\midrule
\textbf{d2} & depth 2, \textbf{unseen inside} & $-1.71$ \tiny[-3.60, +0.31] & \textbf{$+2.13$ \tiny[+0.66, +3.60]} & \textbf{$+3.84$ \tiny[+1.90, +5.70]} \\
\textbf{d3} & depth 3, \textbf{unseen inside} & $-0.35$ \tiny[-2.67, +1.86] & \textbf{$+3.26$ \tiny[+0.81, +5.92]} & \textbf{$+3.60$ \tiny[+1.40, +5.81]} \\
\bottomrule
\end{tabular}
\end{table}

\paragraph{The pre-registered verdict is \textsc{undecided}.} The mapping required \emph{both} that
anchoring hold \emph{and} that breadth fall below the clean-code control; only the first is
established, because the pooled interval at the critical level spans zero. We report that verdict
first.

\paragraph{What is established is stronger than ``no cost''.} On stacks containing an unseen family
the anchored arm is significantly above the clean-code control \emph{and} significantly above
breadth. Read against seen stacks, where it merely matches breadth, the statement is that
\textbf{anchoring's advantage appears only once the input diverges from what breadth was trained
on}.

\paragraph{And breadth's collapse is not uniform.} The pooled figure averages significantly negative
cells with a significantly positive one, split by identifier content (Table~\ref{tab:split}).

\begin{table}[t]\centering\small
\caption{The divergence result across models, at the critical level. \pending\ marks a read that
does not exist yet.}
\label{tab:divmodels}
% GENERATED by scripts/paper/gen_tables.py on 2026-09-12 17:55 UTC
% SOURCE: results/analysis/pipeline/f2_divergence{,_core}_*.json
% Do not edit by hand -- regenerate. Numbers come from result files, never from a report.
\begin{tabular}{llccc}
\toprule
model & lineage & \textsc{cons}$-$\textsc{mono} & \textsc{cons}$-$\textsc{clean} & \textsc{fam}$-$\textsc{clean} \\
\midrule
CodeLlama-7B & Meta & \textbf{$+3.84$ \tiny[+1.90, +5.70]} & \textbf{$+2.13$ \tiny[+0.66, +3.60]} & \textbf{$+5.93$ \tiny[+4.10, +7.78]} \\
CodeLlama-13B & Meta & \textbf{$+3.22$ \tiny[+1.74, +4.65]} & \textbf{$+2.02$ \tiny[+0.58, +3.53]} & \textbf{$+5.97$ \tiny[+4.11, +7.91]} \\
CodeLlama-34B & Meta & \textbf{$+3.95$ \tiny[+2.13, +5.81]} & \textbf{$+3.10$ \tiny[+1.63, +4.58]} & \textbf{$+4.19$ \tiny[+2.48, +5.94]} \\
Llama-3.1-8B & Meta & \textbf{$+2.87$ \tiny[+1.20, +4.62]} & \textbf{$+1.86$ \tiny[+0.27, +3.53]} & \textbf{$+6.67$ \tiny[+4.58, +8.71]} \\
StarCoder2-15B & BigCode & $+1.36$ \tiny[-0.08, +2.71] & $+0.35$ \tiny[-1.13, +1.86] & \textbf{$+3.64$ \tiny[+1.90, +5.35]} \\
Gemma-3-12B & Google & \textbf{$+2.44$ \tiny[+0.97, +3.99]} & $-0.08$ \tiny[-1.67, +1.51] & \textbf{$+5.35$ \tiny[+3.21, +7.37]} \\
CodeGemma-7B & Google & $-1.59$ \tiny[-3.26, +0.12] & \textbf{$-2.40$ \tiny[-4.11, -0.85]} & \textbf{$+5.23$ \tiny[+3.25, +7.21]} \\
Granite-3.1-8B & IBM & $-1.47$ \tiny[-2.99, +0.00] & $-1.86$ \tiny[-3.80, +0.12] & \textbf{$+7.29$ \tiny[+5.19, +9.65]} \\
\bottomrule
\end{tabular}
\end{table}

\paragraph{Scope, in the sentence rather than a footnote.} \slot{final cross-model statement, once
the remaining divergence reads land}

% =======================================================================================
\section{Threats to validity}
\slot{from docs/PAPER_DRAFT.md 9}

\paragraph{Instrument failures we report rather than omit.} The causal attention instrument is
inert: an identifier knockout moves gold log-probability by under 0.5\% for every arm on every
condition, and both of its pre-registered hypotheses were refuted. The attention account is
therefore correlational and is labelled as such.

% =======================================================================================
\section{Related work}
\slot{from docs/PAPER_DRAFT.md 10}

\section{Conclusion}
\slot{once RQ4's cross-model statement is final}

\end{document}
